\section{Related Work}
\label{sec:related}

\subsection{Access control models and context adaptation}

Access control has a long line of models that separate policy from mechanism.
Role-based access control assigns permissions to roles \cite{sandhu1996role};
attribute-based access control generalises the decision to subject, resource and
environment attributes \cite{hu2014guide}; usage control adds continuity and
revocation during use \cite{park2004ucon}. These models provide the vocabulary
for \emph{what} is permitted, but they leave the strength of the enforcing
mechanism to configuration. Zero-trust architectures make the mechanism
explicitly dynamic and assume no implicit trust from network location
\cite{rose2020zero,federici2023zero,zanasi2024flexible,stodt2025beyond}; our work
shares the ``never trust the context unconditionally'' premise but focuses on
the safety of the resulting action rather than on identity verification alone.

Context-adaptive access control (CAAC) makes the posture a function of device
and environmental context \cite{atlam2017overview,rath2017adaptive,genc2019context,arfaoui2019context}.
Recent designs add blockchain-backed continuous authentication
\cite{hussainalnaji2022cab}, smart-contract policy evaluation
\cite{merlec2024sc}, role--capability combinations \cite{krishnasrija2025dynamic},
and app-independent policy enforcement for IoT platforms \cite{mazhar2023maverick}.
Surveys document the breadth of the area and the recurring use of a point risk
estimate as the interface between monitoring and enforcement
\cite{ravidas2019access,qiu2020survey,elhajj2019survey}. Learning-based access
control \cite{nobi2022deep} and context-aware computing more generally
\cite{piccialli2021context} push in the same direction: richer context, more
adaptive decisions. What is usually missing is an explicit adversary who controls
the context channel and an explicit statement of which parts of the decision are
not tradeable. Our contribution is not a new access-control model but a
constrained formulation in which the adaptation is confined by construction to an
already-safe set.

\subsection{Risk estimation and uncertainty}

Risk-based access control uses a score to choose among authentication factors
\cite{lee2017situational,moura2017sotraace}. The score is typically a point
estimate, sometimes with a threshold and sometimes with a cost model. Our belief
is a forward hidden Markov model \cite{rabiner1989tutorial} over latent threat
states, which is a standard partially observable formulation
\cite{kaelbling1998planning}; the difference is that we evaluate it with proper
scoring rules \cite{gneiting2007strictly} and calibration error
\cite{guo2017calibration}, and we do not let it override the invariants.
Self-explainable cyber-physical systems \cite{blumreiter2019self} and
dependability taxonomies \cite{avizienis2004basic} provide complementary
vocabularies for reasoning about why a decision was taken and how it can fail;
we use the latter implicitly in the per-operation failure semantics.

\subsection{Runtime assurance and safety shields}

The closest architectural relatives are runtime assurance and shield synthesis.
The Simplex architecture keeps a verified baseline controller and switches to it
when the advanced controller is untrusted \cite{sha2001simplicity}; safety guards
and shield synthesis enforce a temporal property at the actuator boundary
\cite{wu2017safety,wu2019shield}; the black-box Simplex architecture extends the
idea to learned components \cite{mehmood2022black}; and recent work generates
verified monitors from timed automata \cite{cho2025timed}. Our admissible-set
constructor plays the role of the baseline, and the I4 fallback plays the role of
the switch. The difference is the plant: the actions are access-control postures
in a $243$-element lattice rather than a scalar control input, and the safety
property is a per-operation envelope over authentication, cryptography,
inspection, rate limiting and freshness. The invariant-dominance and
structural-safety propositions are the analogue of the Simplex safety argument
for this lattice. A Simplex-style baseline here would hold a fixed strong posture
as the verified fallback and switch to the adaptive posture only when a monitor
certifies the context; without such a certificate it reduces to Static-High,
which is exactly the low-utility, high-false-denial point in our results. We do
not implement that baseline because the certificate would be the same belief we
are testing, so it would not provide an independent control.

\subsection{Poisoning, Byzantine context and denial of sleep}

Sensor spoofing and false data injection are well studied in cyber-physical
systems \cite{barua2020special,li2019optimal}, and consistency checks can detect
attacks that violate physical invariants \cite{ye2020summation}. Byzantine fault
tolerance bounds the influence of faulty nodes when fewer than half collude
\cite{lamport1982byzantine}, and resilient architectures apply the bound to
industrial and aviation systems \cite{alsulami2021resilient}. Our Byzantine
scenario uses the same $f<N/2$ assumption and adds a cross-node dispersion
signal, but the bound is used only to keep the aggregate informative; it does not
by itself make the interpretation correct. Denial-of-sleep attacks exploit the
reverse incentive, forcing a device to spend energy on security work
\cite{brownfield2005wireless,raymond2009effects}. The energy-inflation and
resonant scenarios are direct stress tests of that failure mode, and the
asymmetric trust update is designed so that a false alarm is corrected slowly, not
immediately amplified.

\subsection{Datasets, calibration and lightweight enforcement}

Public IoT intrusion datasets include IoT-23 \cite{garcia2020iot}, Edge-IIoTset
\cite{ferrag2022edge}, CICIoT2023 \cite{neto2023ciciot2023}, BoT-IoT
\cite{koroniotis2019development}, ToN-IoT \cite{moustafa2020data} and NetFlow-based
evaluation platforms \cite{moustafa2021new}; transformer-based detectors are
commonly evaluated on them \cite{tseng2024multi}. These datasets support
signature and anomaly detection, not access-control decisions under an adaptive
context adversary, which is why we use one of them only to calibrate a prior and
generate the rest of the evaluation from an explicit model. At the enforcement
layer, Ascon is standardised for constrained devices
\cite{dobraunig2021ascon}, and implementation studies report microcontroller
costs \cite{cagua2025implementation}. We treat the cryptographic primitives as
correct and model their cost; the contribution is the decision layer that
selects among them.

\subsection{Positioning}

\Cref{tab:positioning} summarises the distinction. Existing CAAC work optimises
an expected objective against a point estimate and has no non-negotiable floor;
runtime assurance enforces a floor but for a scalar control input rather than an
access-control lattice; Byzantine and poisoning work characterises what an
adversary can do to an aggregate but not what the enforcement point should do
when the aggregate is wrong. This paper combines the three: a constrained CAAC
formulation, an explicit per-operation fail-safe semantics, and an evaluation
that includes the case where the belief is successfully fooled.

\begin{table}[t]
\centering
\small
\caption{Positioning against the three closest lines of work.}
\label{tab:positioning}
\begin{tabular}{p{0.22\linewidth}p{0.30\linewidth}p{0.30\linewidth}}
\toprule
Line of work & What it provides & What it leaves open \\
\midrule
CAAC / risk-based AC
& Context-dependent posture, point risk score
& No non-negotiable floor; estimator error becomes policy error \\
Runtime assurance / shields
& Verified baseline, switching argument
& Scalar control input, not an access-control action lattice \\
Byzantine / poisoning
& Bounds on aggregate influence, attack detectability
& What the enforcement point does when the aggregate is wrong \\
\textbf{This paper}
& Constrained CAAC with I1--I4, per-operation fail-safe action
& Hardware validation, optimal adversary, credential misuse \\
\bottomrule
\end{tabular}
\end{table}
